%% Drafted from the mapped corpus by scripts/draft_topics.py.
% Model: gpt-5.6-terra; source characters: 275279.
\begin{konspekt}
\kreq{Дефиниции (синтаксис)}{език на предикатното смятане от първи ред, терм, формула, област на действие на квантор, свободна и свързана променлива, затворена формула --- \kref{21.1}.}
\kreq{Дефиниции (семантика)}{структура и оценка на индивидните променливи --- \kref{21.2}.}
\kreq{Дефиниции (семантика)}{вярност на формула в структура при оценка, тъждествена вярност, предикатна тавтология --- \kref{21.3}.}
\kreq{Доказателство}{верността при оценка не зависи от стойностите на променливите, които не участват свободно във формулата --- \kref{21.3}.}
\kreq{Формулировка}{лема за субституциите в \emph{общия} й вид --- за произволна формула и допустима субституция --- \kref{21.4}.}
\kreq{Доказателство}{лемата за субституциите в \emph{частния случай} на безкванторна формула. Само този случай се иска с доказателство --- \kref{21.4}.}
\kreq{Типове задачи}{определимост на свойства в дадена структура; изпълнимост на множество формули чрез посочване на модел --- \kref{21.5}.}
\kreq{Литература}{[13], [19], [37].}
\end{konspekt}

\section{Език и синтаксис}\label{s:21.1}

\subsection{Предикатен език от първи ред}\label{s:21.1.1}

\begin{defn}
Езикът на предикатното смятане от първи ред $\mathcal L$ има логическа и нелогическа част.

Логическата част съдържа изброимо множество от индивидуални променливи
\[
\mathrm{Var}=\{x,y,z,x_0,x_1,\ldots\},
\]
логическите връзки $\neg,\wedge,\vee,\Rightarrow,\Leftrightarrow$, кванторите $\forall,\exists$ и помощните символи $(,)$ и запетая. По избор езикът може да съдържа и формален символ за равенство $\doteq$.

Нелогическата част се състои от:
\begin{itemize}
\item множество $\mathrm{Const}_{\mathcal L}$ от индивидуални константи;
\item множество $\mathrm{Func}_{\mathcal L}$ от функционални символи;
\item множество $\mathrm{Pred}_{\mathcal L}$ от предикатни символи.
\end{itemize}
На всеки функционален и предикатен символ се съпоставя положително естествено число, наречено \emph{арност}. Ако $f$ е $n$-местен функционален символ, пишем $\#(f)=n$; аналогично за предикатен символ $p$.
\end{defn}

\begin{remark}
Езикът е от \emph{първи ред}, защото кванторите се отнасят само за елементи на универсума. Не се квантифицира пряко по функции, предикати или множества от елементи.
\end{remark}

\subsection{Термове}\label{s:21.1.2}

\begin{defn}
Множеството на термовете на $\mathcal L$ се дефинира индуктивно:
\begin{enumerate}
\item всяка индивидуална променлива е терм;
\item всяка индивидуална константа е терм;
\item ако $f\in\mathrm{Func}_{\mathcal L}$, $\#(f)=n$, и $\tau_1,\ldots,\tau_n$ са термове, то
\[
f(\tau_1,\ldots,\tau_n)
\]
е терм.
\end{enumerate}
\end{defn}

\begin{example}
Ако $c$ е константа, $f$ е едноместен, а $g$ е двуместен функционален символ, то
\[
x,\qquad c,\qquad f(x),\qquad g(f(x),c)
\]
са термове.
\end{example}

За терм $\tau$ с $\operatorname{Var}[\tau]$ означаваме множеството от променливите, които участват в него:
\[
\operatorname{Var}[x]=\{x\},\qquad
\operatorname{Var}[c]=\varnothing,
\]
\[
\operatorname{Var}[f(\tau_1,\ldots,\tau_n)]
 =\bigcup_{i=1}^{n}\operatorname{Var}[\tau_i].
\]
\begin{defn}[Затворен терм]
Терм без променливи се нарича \emph{затворен} или \emph{основен} терм.
\end{defn}

\subsection{Формули}\label{s:21.1.3}

\begin{defn}
Ако $p\in\mathrm{Pred}_{\mathcal L}$ е $n$-местен и $\tau_1,\ldots,\tau_n$ са термове, то
\[
p(\tau_1,\ldots,\tau_n)
\]
се нарича \emph{атомарна формула}. Ако езикът е с формално равенство, то и
\[
\tau_1\doteq\tau_2
\]
е атомарна формула.
\end{defn}

\begin{defn}
Множеството на формулите на $\mathcal L$ се задава индуктивно:
\begin{enumerate}
\item атомарните формули са формули;
\item ако $\varphi$ и $\psi$ са формули, то $\neg\varphi$,
\[
(\varphi\wedge\psi),\quad
(\varphi\vee\psi),\quad
(\varphi\Rightarrow\psi),\quad
(\varphi\Leftrightarrow\psi)
\]
са формули;
\item ако $x$ е индивидуална променлива и $\varphi$ е формула, то
\[
\forall x\,\varphi,\qquad \exists x\,\varphi
\]
са формули.
\end{enumerate}
\end{defn}

\begin{defn}
Формула $\psi$ е \emph{подформула} на формула $\varphi$, ако $\psi$ се получава като част от синтактичното построение на $\varphi$.
\end{defn}

\subsection{Област на действие; свободни и свързани променливи}\label{s:21.1.4}

\begin{defn}
В подформула от вида
\[
\forall x\,\psi\qquad\text{или}\qquad\exists x\,\psi
\]
формулата $\psi$ се нарича \emph{област на действие} на съответния квантор.
\end{defn}

\begin{defn}
Едно участие на променливата $x$ във формула е \emph{свързано}, ако лежи в областта на действие на квантор по $x$. В противен случай участието е \emph{свободно}.
\end{defn}

Най-близкият квантор по $x$, чиято област на действие съдържа дадено участие, определя дали това участие е свързано. Една и съща променлива може да има както свободни, така и свързани участия.

\begin{example}
Във формулата
\[
p(x)\wedge\forall x\,q(x,y)
\]
първото участие на $x$ е свободно, а второто е свързано. Всички участия на $y$ са свободни.
\end{example}

С $\operatorname{Var}^{\mathrm{free}}[\varphi]$ означаваме множеството на свободните променливи на $\varphi$, а с $\operatorname{Var}^{\mathrm{bd}}[\varphi]$ --- множеството на променливите с поне едно свързано участие. Важни рекурсивни равенства са:
\[
\operatorname{Var}^{\mathrm{free}}[p(\tau_1,\ldots,\tau_n)]
=\bigcup_{i=1}^{n}\operatorname{Var}[\tau_i],
\]
\[
\operatorname{Var}^{\mathrm{free}}[\neg\varphi]
=\operatorname{Var}^{\mathrm{free}}[\varphi],
\]
\[
\operatorname{Var}^{\mathrm{free}}[\varphi\circ\psi]
=\operatorname{Var}^{\mathrm{free}}[\varphi]\cup
\operatorname{Var}^{\mathrm{free}}[\psi],
\qquad
\circ\in\{\wedge,\vee,\Rightarrow,\Leftrightarrow\},
\]
\[
\operatorname{Var}^{\mathrm{free}}[\forall x\,\varphi]
=\operatorname{Var}^{\mathrm{free}}[\varphi]\setminus\{x\},
\qquad
\operatorname{Var}^{\mathrm{free}}[\exists x\,\varphi]
=\operatorname{Var}^{\mathrm{free}}[\varphi]\setminus\{x\}.
\]

\begin{defn}
Формула $\varphi$ се нарича \emph{затворена}, ако
\[
\operatorname{Var}^{\mathrm{free}}[\varphi]=\varnothing.
\]
\end{defn}

\section{Структури, оценки и стойности на термове}\label{s:21.2}

\begin{defn}
Структура за езика $\mathcal L$ е двойка
\[
\mathcal A=\langle A,I\rangle,
\]
където $A\neq\varnothing$ е \emph{универсум} или \emph{носител} на структурата, а $I$ е интерпретация на нелогическите символи:
\begin{itemize}
\item за всяка константа $c$ е зададен елемент $c^{\mathcal A}\in A$;
\item за всеки $n$-местен функционален символ $f$ е зададена функция
\[
f^{\mathcal A}:A^n\to A;
\]
\item за всеки $n$-местен предикатен символ $p$ е зададена релация
\[
p^{\mathcal A}\subseteq A^n.
\]
\end{itemize}
\end{defn}

\begin{remark}
Функционалният символ $f$ е синтактичен знак, докато $f^{\mathcal A}$ е функция върху универсумa. Аналогично $p$ е знак, а $p^{\mathcal A}$ е релация.
\end{remark}

\begin{example}
За език с константа $0$, двуместен функционален символ $+$ и двуместен предикатен символ $<$ естествените числа могат да се разглеждат като структура
\[
\mathcal N=\langle\mathbb N,0^{\mathcal N},+^{\mathcal N},<^{\mathcal N}\rangle.
\]
\end{example}

\begin{defn}
\emph{Оценка на индивидуалните променливи} в структура $\mathcal A$ е функция
\[
v:\mathrm{Var}\to A.
\]
За $a\in A$ с $v_x^a$ означаваме оценката, модифицирана в $x$ с $a$:
\[
v_x^a(y)=
\begin{cases}
a,&y=x,\\
v(y),&y\ne x.
\end{cases}
\]
\end{defn}

Стойността на терм $\tau$ в $\mathcal A$ при оценка $v$, означавана с
\[
\llbracket\tau\rrbracket^{\mathcal A}_v,
\]
се определя рекурсивно:
\[
\llbracket x\rrbracket^{\mathcal A}_v=v(x),
\qquad
\llbracket c\rrbracket^{\mathcal A}_v=c^{\mathcal A},
\]
\[
\llbracket f(\tau_1,\ldots,\tau_n)\rrbracket^{\mathcal A}_v
=
f^{\mathcal A}\bigl(
\llbracket\tau_1\rrbracket^{\mathcal A}_v,\ldots,
\llbracket\tau_n\rrbracket^{\mathcal A}_v
\bigr).
\]

\begin{prop}
Нека $\tau$ е терм, а $v_1,v_2$ са оценки в $\mathcal A$. Ако
\[
v_1(x)=v_2(x)
\quad\text{за всяко }x\in\operatorname{Var}[\tau],
\]
то
\[
\llbracket\tau\rrbracket^{\mathcal A}_{v_1}
=
\llbracket\tau\rrbracket^{\mathcal A}_{v_2}.
\]
\end{prop}

\begin{proof}
Доказателството е индукция по построението на $\tau$. За променлива твърдението следва от условието; за константа стойността не зависи от оценката. При
\[
\tau=f(\tau_1,\ldots,\tau_n)
\]
оценките съвпадат върху променливите на всеки $\tau_i$. Индукционната хипотеза дава равенство на стойностите на всички аргументи, а после и равенство на стойностите след прилагане на $f^{\mathcal A}$.
\end{proof}

\begin{cor}
Стойността на затворен терм не зависи от оценката.
\end{cor}

\begin{proof}
Множеството от променливи на затворен терм е празно. Всякакви две оценки съвпадат върху това множество. Предходното твърдение за зависимостта само от променливите на терма тогава дава еднакви стойности при двете оценки.
\end{proof}

\section{Семантика на формулите}\label{s:21.3}

\subsection{Вярност при оценка}\label{s:21.3.1}

За формула $\varphi$, структура $\mathcal A$ и оценка $v$ записваме
\[
\mathcal A\models\varphi[v]
\]
и четем: „$\varphi$ е вярна в $\mathcal A$ при оценка $v$“. Релацията се определя рекурсивно.

За атомарна формула:
\[
\mathcal A\models p(\tau_1,\ldots,\tau_n)[v]
\quad\Longleftrightarrow\quad
\bigl(
\llbracket\tau_1\rrbracket^{\mathcal A}_v,\ldots,
\llbracket\tau_n\rrbracket^{\mathcal A}_v
\bigr)\in p^{\mathcal A}.
\]
При наличие на формално равенство:
\[
\mathcal A\models(\tau_1\doteq\tau_2)[v]
\quad\Longleftrightarrow\quad
\llbracket\tau_1\rrbracket^{\mathcal A}_v
=
\llbracket\tau_2\rrbracket^{\mathcal A}_v.
\]

За логическите връзки се използва класическата двузначна семантика:
\[
\mathcal A\models\neg\varphi[v]
\quad\Longleftrightarrow\quad
\mathcal A\not\models\varphi[v],
\]
\[
\mathcal A\models(\varphi\wedge\psi)[v]
\quad\Longleftrightarrow\quad
\mathcal A\models\varphi[v]\ \text{и}\
\mathcal A\models\psi[v],
\]
\[
\mathcal A\models(\varphi\vee\psi)[v]
\quad\Longleftrightarrow\quad
\mathcal A\models\varphi[v]\ \text{или}\
\mathcal A\models\psi[v],
\]
\[
\mathcal A\models(\varphi\Rightarrow\psi)[v]
\quad\Longleftrightarrow\quad
\mathcal A\not\models\varphi[v]\ \text{или}\
\mathcal A\models\psi[v].
\]

За кванторите:
\[
\mathcal A\models\forall x\,\varphi[v]
\quad\Longleftrightarrow\quad
\text{за всяко }a\in A:
\mathcal A\models\varphi[v_x^a],
\]
\[
\mathcal A\models\exists x\,\varphi[v]
\quad\Longleftrightarrow\quad
\text{съществува }a\in A:
\mathcal A\models\varphi[v_x^a].
\]

\begin{example}
Нека $p^{\mathcal A}\subseteq A$ е едноместна релация. Тогава
\[
\mathcal A\models\forall x\,p(x)[v]
\]
означава, че $p^{\mathcal A}=A$, а
\[
\mathcal A\models\exists x\,p(x)[v]
\]
означава, че $p^{\mathcal A}\ne\varnothing$.
\end{example}

\subsection{Зависимост само от свободните променливи}\label{s:21.3.2}

\begin{keythm}
Нека $\varphi$ е формула, а $v_1,v_2$ са оценки в структурата $\mathcal A$. Ако
\[
v_1(x)=v_2(x)
\quad\text{за всяко }
x\in\operatorname{Var}^{\mathrm{free}}[\varphi],
\]
то
\[
\mathcal A\models\varphi[v_1]
\quad\Longleftrightarrow\quad
\mathcal A\models\varphi[v_2].
\]
\end{keythm}

\begin{proof}
Доказваме с индукция по построението на $\varphi$.

Нека $\varphi=p(\tau_1,\ldots,\tau_n)$. Всички променливи на $\tau_i$ са свободни във формулата. По твърдението за термовете
\[
\llbracket\tau_i\rrbracket^{\mathcal A}_{v_1}
=
\llbracket\tau_i\rrbracket^{\mathcal A}_{v_2}
\qquad (1\leq i\leq n).
\]
Следователно едната наредена $n$-торка принадлежи на $p^{\mathcal A}$ точно тогава, когато принадлежи другата. Случаят с формално равенство е аналогичен.

При отрицание и при двоичните логически връзки твърдението следва непосредствено от индукционната хипотеза и класическата семантика.

Нека
\[
\varphi=\forall x\,\psi.
\]
Тогава свободните променливи на $\varphi$ са свободните променливи на $\psi$, различни от $x$. За произволно $a\in A$ оценките $v_{1\,x}^a$ и $v_{2\,x}^a$ съвпадат върху всички свободни променливи на $\psi$: за $x$ и двете дават $a$, а за останалите следва от условието. По индукционната хипотеза
\[
\mathcal A\models\psi[v_{1\,x}^a]
\Longleftrightarrow
\mathcal A\models\psi[v_{2\,x}^a].
\]
Следователно всеобщото квантифициране дава еднаква стойност. Случаят $\exists x\,\psi$ е аналогичен.
\end{proof}

\begin{cor}
Ако $\varphi$ е затворена формула, нейната вярност в дадена структура не зависи от оценката.
\end{cor}

\begin{proof}
За затворена формула $\operatorname{FV}(\varphi)=\varnothing$. Всякакви две оценки съвпадат върху свободните й променливи. По доказаната теорема за независимост от останалите променливи истинностната стойност е една и съща.
\end{proof}

\subsection{Тъждествена вярност и предикатни тавтологии}\label{s:21.3.3}

\begin{defn}
Формулата $\varphi$ е \emph{тъждествено вярна в структурата} $\mathcal A$, ако
\[
\mathcal A\models\varphi[v]
\]
за всяка оценка $v$ в $\mathcal A$. Пишем
\[
\mathcal A\models\varphi.
\]
\end{defn}

\begin{defn}
Формулата $\varphi$ е \emph{предикатна тавтология} или \emph{общовалидна}, ако е тъждествено вярна във всяка структура за езика:
\[
\models\varphi.
\]
\end{defn}

\begin{remark}
За затворена формула е достатъчно да се провери една произволна оценка: по предходното следствие всички оценки дават една и съща истинностна стойност.
\end{remark}

\section{Субституции}\label{s:21.4}

\subsection{Субституция на термове}\label{s:21.4.1}

\begin{defn}
Субституция е съпоставяне, което на всяка индивидуална променлива задава терм. Често се задава крайно чрез
\[
s=\left\{\frac{x_1}{t_1},\ldots,\frac{x_n}{t_n}\right\},
\]
където $x_1,\ldots,x_n$ са различни променливи. За терм $\tau$ с $\tau s$ означаваме терма, получен при едновременно заместване на съответните променливи.
\end{defn}

\begin{lem}[Лема за субституциите на термове]
Нека $\tau(x_1,\ldots,x_n)$ е терм, $t_1,\ldots,t_n$ са термове, а $v$ е оценка в $\mathcal A$. Ако $\omega$ е оценка, зададена чрез
\[
\omega(x_i)=\llbracket t_i\rrbracket^{\mathcal A}_v
\quad (1\leq i\leq n),
\]
и съвпадаща с $v$ върху останалите променливи, то
\[
\llbracket\tau[x_1/t_1,\ldots,x_n/t_n]\rrbracket^{\mathcal A}_v
=
\llbracket\tau\rrbracket^{\mathcal A}_\omega.
\]
\end{lem}

\begin{proof}
Доказателството е индукция по построението на $\tau$. За $\tau=x_i$ равенството следва от дефиницията на $\omega$; за променлива, която не се заменя, и за константа е непосредствено. При
\[
\tau=f(\tau_1,\ldots,\tau_k)
\]
се прилага индукционната хипотеза към всеки аргумент и после дефиницията на стойността на функционален терм.
\end{proof}

\subsection{Допустимост на заместването във формули}\label{s:21.4.2}

При замяна във формула се заменят само свободните участия на съответната променлива. Трябва да се избягва \emph{захващане} на променлива.

\begin{defn}
Променлива $y$ е опасна за формула $\varphi$ при субституция $s$, ако $y$ участва в терма $s(x)$ за някоя свободна променлива $x$ на $\varphi$, за която $s(x)\ne x$.
\end{defn}

Заместването е допустимо, когато свободно участие, заменено с терм, не попада в област на действие на квантор по променлива, участваща в този терм.

\begin{example}
Нека
\[
\varphi=\forall y\,(q(x,y)\Rightarrow p(y)).
\]
Замяната на свободното $x$ с $f(y)$ не е допустима: полученото
\[
\forall y\,(q(f(y),y)\Rightarrow p(y))
\]
променя ролята на $y$, защото то става свързано от квантора.
\end{example}

\begin{defn}[Допустима субституция]
Субституцията $s$ е \emph{допустима} за формулата $\varphi$, ако за всяка
свободна променлива $x$ на $\varphi$ с $s(x)\ne x$ никое свободно участие на $x$
в $\varphi$ не попада в областта на действие на квантор по променлива, която
участва в терма $s(x)$. Тоест никоя променлива на заместващия терм не се
свързва при заместването.
\end{defn}

Конспектът иска твърдението за стойността при замяна на променливи с термове да
бъде \emph{формулирано} в общия си вид и да бъде \emph{доказано} в частния случай
на безкванторна формула. Затова тук стоят и двете: общата формулировка, която
трябва да може да се възпроизведе точно, и подробната индукция за безкванторния
случай.

\begin{keythm}[Лема за субституциите, общ вид]
Нека $\varphi$ е произволна формула, $s$ е субституция, \emph{допустима} за
$\varphi$, $v$ е оценка в структурата $\mathcal A$, а $\omega$ е оценката
\[
\omega(x)=\llbracket xs\rrbracket^{\mathcal A}_v
\]
за всяка променлива $x$. Тогава
\[
\mathcal A\models(\varphi s)[v]
\quad\Longleftrightarrow\quad
\mathcal A\models\varphi[\omega].
\]
\end{keythm}

\begin{remark}
Изискването за допустимост не може да се пропусне. За
$\varphi=\exists y\,(x\ne y)$ и $s(x)=y$ лявата страна дава
$\exists y\,(y\ne y)$, което е невярно във всяка структура, докато дясната е
вярна във всяка структура с поне два елемента. Точно това е захващането на
променлива от предходния пример.
\end{remark}

\begin{proof}
Индукция по построението на $\varphi$. Атомарните случаи и случаите с
$\neg,\wedge,\vee,\Rightarrow,\Leftrightarrow$ не използват допустимостта и са
изписани подробно в теоремата за безкванторния случай по-долу; повторени тук,
те дават базата и пропозиционалните стъпки.

Остава случаят с квантор. Нека $\varphi=\forall y\,\psi$ и нека $s_y$ е
субституцията, която съвпада с $s$ извън $y$ и полага $s_y(y)=y$: заместват се
само свободните участия, а $y$ е свързана в $\varphi$. Тогава
$\varphi s=\forall y\,(\psi s_y)$ и $s_y$ е допустима за $\psi$. По дефиниция на
семантиката
\[
\mathcal A\models(\varphi s)[v]
\quad\Longleftrightarrow\quad
\text{за всяко }a\in A:\ \mathcal A\models(\psi s_y)[v_y^a].
\]
Прилагаме индукционната хипотеза за $\psi$, $s_y$ и оценката $v_y^a$. Тя дава
\[
\mathcal A\models(\psi s_y)[v_y^a]
\quad\Longleftrightarrow\quad
\mathcal A\models\psi[\omega_a],
\qquad
\omega_a(x)=\llbracket xs_y\rrbracket^{\mathcal A}_{v_y^a}.
\]
Сравняваме $\omega_a$ с $\omega_y^a$ върху свободните променливи на $\psi$.
За $x=y$ имаме $ys_y=y$, тоест $\omega_a(y)=a=\omega_y^a(y)$. За свободна
променлива $x\ne y$ на $\psi$, която е свободна и в $\varphi$, имаме
$xs_y=xs$. Ако $s(x)=x$, стойността е $v(x)$ и в двата случая. Ако
$s(x)\ne x$, допустимостта на $s$ за $\varphi$ казва, че $y$ не участва в
терма $s(x)$; тогава по твърдението за зависимостта на стойността на терм само
от участващите в него променливи
\[
\llbracket s(x)\rrbracket^{\mathcal A}_{v_y^a}
=
\llbracket s(x)\rrbracket^{\mathcal A}_{v}
=\omega(x)=\omega_y^a(x).
\]
Следователно $\omega_a$ и $\omega_y^a$ съвпадат върху
$\operatorname{Var}^{\mathrm{free}}[\psi]$ и по теоремата за зависимост само от
свободните променливи
\[
\mathcal A\models\psi[\omega_a]
\quad\Longleftrightarrow\quad
\mathcal A\models\psi[\omega_y^a].
\]
Събирайки еквивалентностите за всяко $a\in A$, получаваме
\[
\mathcal A\models(\varphi s)[v]
\ \Longleftrightarrow\
\text{за всяко }a\in A:\ \mathcal A\models\psi[\omega_y^a]
\ \Longleftrightarrow\
\mathcal A\models\forall y\,\psi\,[\omega].
\]
Случаят $\varphi=\exists y\,\psi$ се получава по същия начин, като „за всяко
$a$“ се замени със „съществува $a$“.
\end{proof}

\begin{keythm}[Лема за субституциите, безкванторен случай]
Нека $\varphi$ е безкванторна формула, $s$ е субституция, $v$ е оценка в $\mathcal A$, а $\omega$ е оценката
\[
\omega(x)=\llbracket xs\rrbracket^{\mathcal A}_v
\]
за всяка променлива $x$. Тогава
\[
\mathcal A\models(\varphi s)[v]
\quad\Longleftrightarrow\quad
\mathcal A\models\varphi[\omega].
\]
\end{keythm}

Това е \emph{частният случай}, чието доказателство конспектът изисква. При
безкванторна $\varphi$ всяка субституция е допустима, защото няма квантор, който
да свърже променлива, затова хипотезата за допустимост отпада.

\begin{proof}
Доказваме с индукция по построението на безкванторната формула $\varphi$.

Нека
\[
\varphi=p(\tau_1,\ldots,\tau_n).
\]
По лемата за субституциите на термове:
\[
\llbracket\tau_i s\rrbracket^{\mathcal A}_v
=
\llbracket\tau_i\rrbracket^{\mathcal A}_\omega
\qquad (1\leq i\leq n).
\]
Следователно
\[
\begin{aligned}
\mathcal A\models(\varphi s)[v]
&\Longleftrightarrow
\bigl(
\llbracket\tau_1s\rrbracket^{\mathcal A}_v,\ldots,
\llbracket\tau_ns\rrbracket^{\mathcal A}_v
\bigr)\in p^{\mathcal A}\\
&\Longleftrightarrow
\bigl(
\llbracket\tau_1\rrbracket^{\mathcal A}_\omega,\ldots,
\llbracket\tau_n\rrbracket^{\mathcal A}_\omega
\bigr)\in p^{\mathcal A}\\
&\Longleftrightarrow
\mathcal A\models\varphi[\omega].
\end{aligned}
\]
Случаят с равенство е аналогичен.

Нека $\varphi=\neg\psi$. По индукционната хипотеза
\[
\mathcal A\models(\psi s)[v]
\Longleftrightarrow
\mathcal A\models\psi[\omega].
\]
След отрицание получаваме търсената еквивалентност за $\neg\psi$.

Нека
\[
\varphi=(\psi\wedge\chi).
\]
Тогава
\[
(\psi\wedge\chi)s=(\psi s\wedge\chi s).
\]
По индукционната хипотеза:
\[
\mathcal A\models(\psi s)[v]\Longleftrightarrow\mathcal A\models\psi[\omega],
\]
\[
\mathcal A\models(\chi s)[v]\Longleftrightarrow\mathcal A\models\chi[\omega].
\]
Следователно
\[
\mathcal A\models((\psi\wedge\chi)s)[v]
\Longleftrightarrow
\mathcal A\models(\psi\wedge\chi)[\omega].
\]
Случаите за $\vee$, $\Rightarrow$ и $\Leftrightarrow$ са аналогични. Понеже $\varphi$ е безкванторна, случаи с квантори няма.
\end{proof}

\section{Определимост и изпълнимост}\label{s:21.5}

\subsection{Определими свойства}\label{s:21.5.1}

Нека $\varphi[x_1,\ldots,x_n]$ означава, че всички свободни променливи на $\varphi$ са сред $x_1,\ldots,x_n$. Формулата определя в структурата $\mathcal A$ множеството
\[
D_\varphi=
\left\{
(a_1,\ldots,a_n)\in A^n
\mid
\mathcal A\models\varphi[v]
\text{ за оценка }v(x_i)=a_i
\right\}.
\]
Това е коректно определено именно защото истинността на формулата зависи само от свободните й променливи.

\begin{example}
В структурата
\[
\mathcal N=\langle\mathbb N,0,+,<\rangle
\]
формулата
\[
\exists z\,(x=z+z)
\]
определя множеството на четните естествени числа.
\end{example}

\subsection{Изпълнимост на множества от формули}\label{s:21.5.2}

\begin{defn}
Множество $\Gamma$ от формули е \emph{изпълнимо в структурата} $\mathcal A$, ако съществува оценка $v$, за която
\[
\mathcal A\models\psi[v]
\qquad\text{за всяка }\psi\in\Gamma.
\]
\end{defn}

\begin{defn}
Структурата $\mathcal A$ е \emph{модел} на $\Gamma$, ако
\[
\mathcal A\models\psi
\qquad\text{за всяка }\psi\in\Gamma.
\]
Пишем
\[
\mathcal A\models\Gamma.
\]
\end{defn}

\begin{example}
Нека езикът съдържа едноместен предикат $P$ и двуместен предикат $R$. Множеството
\[
\Gamma=
\left\{
\forall x\,\exists y\,R(x,y),\
\forall x\,\bigl(P(x)\Rightarrow\exists y\,R(y,x)\bigr),\
\exists x\,P(x)
\right\}
\]
е изпълнимо. Модел е структура с универсум $A=\{a\}$, за която
\[
P^{\mathcal A}=\{a\},
\qquad
R^{\mathcal A}=\{(a,a)\}.
\]
Действително, всяка квантифицирана променлива може да получи стойност $a$, а всички формули от $\Gamma$ са верни.
\end{example}

\begin{supp}[Бележка]
При формули със свободни променливи трябва да се различават:
\[
\text{„има една оценка, която удовлетворява всички формули“}
\]
и
\[
\text{„всички формули са тъждествено верни в структурата“}.
\]
Второто е по-силно условие. За затворени формули двете понятия съвпадат, защото истинността не зависи от оценката.
\end{supp}

\section{Изпитен фокус}\label{s:21.6}

\begin{itemize}
\item Да се дефинират език от първи ред, арност, терм, атомарна формула и предикатна формула.
\item Да се обяснят област на действие на квантор, свободно и свързано участие, свободна променлива и затворена формула.
\item Да се дефинират структура, универсум, интерпретация, оценка и модифицирана оценка $v_x^a$.
\item Да се даде рекурсивната семантика на термовете и формулите, особено на $\forall$ и $\exists$.
\item Да се различават вярност при оценка $\mathcal A\models\varphi[v]$, тъждествена вярност $\mathcal A\models\varphi$ и предикатна тавтология $\models\varphi$.
\item Да се възпроизведе доказателствената идея, че стойността на формула зависи само от свободните й променливи: индукция по построението на формулата; при кванторите се сравняват модифицираните оценки.
\item Да се формулира лемата за субституциите на термове.
\item Да се формулира лемата за субституциите в \emph{общия} й вид --- за произволна формула $\varphi$ и субституция $s$, \emph{допустима} за $\varphi$: $\mathcal A\models(\varphi s)[v]\Leftrightarrow\mathcal A\models\varphi[\omega]$, където $\omega(x)=\llbracket xs\rrbracket^{\mathcal A}_v$. Конспектът иска формулировката в общия вид.
\item Да се докаже лемата за субституциите в частния случай на \emph{безкванторна} формула чрез индукция по построението й. Само този частен случай се изисква с доказателство.
\item Да се обясни защо е нужна допустимост на заместването и да се даде пример за захващане на променлива.
\item Да се решават задачи за определимост чрез формула и за изпълнимост чрез явно посочване на структура, интерпретации и при нужда оценка.
\end{itemize}
